\documentclass[a4paper,10pt]{article}
\usepackage[spanish]{babel}
\usepackage{
			ucs
			,amsfonts
			,amssymb
			,acronym
			,verbatim
			,rotating
			,graphicx
			,enumerate
			,fancyhdr
            ,mathrsfs
            ,array
			}

\usepackage{tikz}
\usetikzlibrary{svg.path}

\pagestyle{fancy}
\addtolength{\textheight}{1.0in}
\addtolength{\textwidth}{0.5in}
\addtolength{\voffset}{-1.0in}

\fancyhf{}
\fancyhead[L]{{\bfseries Programaci\'on I Curso 2012. Pr\'actica 8 \hfill Listas}}
\fancyfoot[LE,RO]{\thepage}

\fancypagestyle{firstpage}{
\fancyhf{}

\fancyhead[L]{\includegraphics[width=1.75cm]{ucalp.jpg}  \\
\small{UCALP. Facultad de Ciencias Exactas e Ingenier\'ia} \\
{\bfseries Programaci\'on I Curso 2012. Pr\'actica 8 \hfill Listas}}
\fancyfoot[LE,RO]{\thepage}
}

\thispagestyle{firstpage}

\begin{document}

~\\ \\ \\ \\

\begin{enumerate}

\item Modifique los algoritmos de ordenamiento y b\'usqueda vistos con anterioridad para que funcionen sobre listas enlazadas.

\item Escribir un procedimiento/funci\'on que permita intercambiar los elementos ubicados en la posici\'on \verb-p- y \verb-siguiente(p)- de una lista enlazada.
     
\item Escribir funciones/procedimientos que permitan \emph{insertar, borrar} y \emph{localizar} un elemento en un lista enlazada, pero de manera ordenada (\emph{ascendente o descendente, es indistinto}).
    
\item Escriba una funci\'on/procedimiento que permita eliminar todas las ocurrencias de un elemento $x$ dentro de una lista enlazada.

\item Escriba una funci\'on/procedimiento que permita ``dar vuelta'' una lista enlazada.

\item Implemente el ejercicio 3 de la pr\'actica anterior (\emph{7.5}) utilizando listas enlazadas.   

\item Implemente el ejercicio 4 de la pr\'actica anterior (\emph{7.5}) utilizando listas enlazadas.

\item Implemente el ejercicio 5 de la pr\'actica anterior (\emph{7.5}) utilizando listas enlazadas.

\item Se desea contar con un tipo \emph{String} de longitud variable (una lista de caracteres). Realice la implementaci\'on de este nuevo tipo y funciones/procedimientos para manejo: \emph{cargar, mostrar, saber longitud, comparar, etc...}

\item Realizar una funci\'on/procedimiento que permita chequear si una palabra (representada como en el punto anterior) es palindrome o no. 
     
\end{enumerate}
\end{document}
